\Needspace{9\baselineskip}
\section{The minimum principle}
\label{sec:psh-minimum}

An infimum of plurisubharmonic functions need not be plurisubharmonic. Kiselman's principle gives a different operation: take an infimum in a real parameter whose complexification is plurisubharmonic and invariant in the imaginary direction. Freezing the background coefficients preserves the precise lower curvature bound.

\begin{lemma}\label{lem:invariant-strip-convexity}
Let $I\subset\mathbb R$ be a nonempty open interval, set $S_I=\{z\in\mathbb C:\Real z\in I\}$, and let $\pi:X\times S_I\to X$ be the projection. Suppose $U(x,z)=u_{\Real z}(x)$ is quasi-plurisubharmonic and $\pi^*\theta+\ddc U\geq0$. For each $x$, the function $t\mapsto u_t(x)$ is finite and convex or identically $-\infty$. The finite alternative holds for almost every $x$.
\end{lemma}
\begin{proof}
A local smooth correction making $U$ plurisubharmonic can be chosen to depend only on $x$. Restriction to a vertical complex line is therefore subharmonic or identically $-\infty$. In the former case its distributional Laplacian is the second real derivative of $t\mapsto u_t(x)$, which is nonnegative. The resulting convex representative is finite on $I$ and agrees everywhere with the subharmonic function by the submean inequality. Finally, $U$ is locally integrable on the product. Fubini's theorem, applied in coordinate balls times a fixed rectangle in $S_I$, excludes the identically $-\infty$ alternative for almost every $x$.
\end{proof}

\begin{proposition}\label{prop:hermitian-kiselman}
Under the hypotheses of \cref{lem:invariant-strip-convexity}, the function
\[
 q(x)=\inf_{t\in I}u_t(x)
\]
is either a genuine $\theta$-plurisubharmonic function or identically $-\infty$. No regularization of the infimum is necessary. If $I$ is bounded, the infimum is genuine.
\end{proposition}
\begin{proof}
The classical local Kiselman principle says that the infimum in the real variable of a plurisubharmonic function on a product with a strip, invariant in its imaginary variable, is plurisubharmonic or identically $-\infty$. Its parameter interval may be unbounded. We use this local result from \cite{SGPT}.

Fix $x_0\in X$, choose local holomorphic coordinates centered at $x_0$, and let $\kappa$ be their Euclidean form. Fix $\eta>0$. On a sufficiently small coordinate ball the constant form $\beta=\theta(x_0)$ satisfies
\[
 \beta-\eta\kappa\leq\theta\leq\beta+\eta\kappa.
\]
Choose a quadratic function $r$ with $\ddc r=\beta+\eta\kappa$. Then $U+r\circ\pi$ is psh on the ball times $S_I$. Its infimum is $q+r$, so the local principle gives
\[
 \ddc q\geq-\beta-\eta\kappa\geq-\theta-2\eta\kappa
\]
where the infimum is genuine. In a fixed atlas the Euclidean forms are bounded by a fixed Hermitian metric. Letting $\eta\downarrow0$ proves the exact lower bound $\ddc q\geq-\theta$. The local alternative propagates across overlapping balls on connected $X$: a genuine psh function cannot equal $-\infty$ on an open set. Thus either $q$ is genuine everywhere as a quasi-plurisubharmonic function or $q\equiv-\infty$. The infimum is upper semicontinuous, as an infimum of upper semicontinuous functions.

Suppose $I$ is bounded. By \cref{lem:invariant-strip-convexity}, the function $t\mapsto u_t(x)$ is finite and convex for almost every $x$. A supporting affine function at any interior point bounds it from below on $I$. Hence $q(x)>-\infty$ almost everywhere, excluding the identically $-\infty$ alternative.
\end{proof}

For example, the proposition applies to $u_t-t\tau$ for a fixed real $\tau$. On an unbounded interval the infimum may be $-\infty$; this alternative will encode the cutoff of a test curve.
